% Physical page 5 onward: the developed commentary of the concise study,
% organized by the questions the three readings answer differently rather than
% by the order of the Mass. Each unit draws on several appointed elements and
% sets the readings' answers beside one another.

\section{The Propers: Detailed Commentary}
\zlabel{triptych:concise:commentary:start}

\subsection{Who was invited, and where the call went}

The first half of the parable is where most of its words are spent, and the
second reading hears the whole Mass from it. St John Chrysostom reads the
parable where St Matthew places it, as the answer that follows the vineyard. At
the end of the parable of the husbandmen Christ had said that the kingdom
``shall be given to a nation bringing forth the fruits thereof'', and the
marriage feast ``declares next to what kind of nation''. It ``proclaims
beforehand both the casting out of the Jews, and the calling of the Gentiles''
(\work{Homilies on Matthew} 69). The invitation had come many times, ``By all
the prophets; by John again \dots\ by the Son Himself again'', and after the
Ascension ``by Peter, and those with him''. The apostles went to the Gentiles
only when they had been driven out, and Chrysostom bids his hearers hear Paul
interpreting the parable: ``It was necessary that the word of God should first have
been spoken to you, but since ye judge yourselves unworthy, lo, we turn to the
Gentiles'' (Acts 13:46).

St Irenaeus used the parable in the second century against teachers who divided
the God of the prophets from the Father of Jesus Christ. It shows ``that there
is one King and Lord, the Father of all'', who called the men of the former
dispensation by his servants and, when they refused, ``called together from all
the highways, that is, from all nations, [guests] to the marriage feast of His
Son''. ``The Lord, therefore, who has called us everywhere by the apostles, is
He who called those of old by the prophets'' (\work{Against Heresies}
IV.36.5). That proof lets the Introit's Lord, who says ``my people'' in the
psalm of Israel's history, be heard as the Lord who calls the Gentiles in the
Alleluia and the Gospel. St Hilary of Poitiers names the first invited
\latin{populus Israel}, called through the law to the glory of eternity; the
servants sent to call them are the apostles, and those sent a second time
apostolic men, the apostles' successors; and because the parable speaks of the
time of judgement and resurrection, the king turns the same word \latin{ad
congregationem gentium}, to the gathering of the nations (\work{Commentarius in
Matthaeum} 22, 3--6). St Jerome agrees, and adds his own details. The wedding
is made for Christ and the Church, \latin{quae tam ex Judaeis, quam ex gentibus,
congregata est}; the first servant is Moses or, reading the plural, the
prophets; and \latin{Gentilium populus non erat in viis, sed in exitibus
viarum}, the people of the Gentiles was not in the ways but at the ends of the
ways (\work{In Matthaeum} III).

St Augustine answers the question of the highways in a form of his own:
\latin{viæ intelliguntur dogmata Gentium}, the ways are the teachings of the
Gentiles, because out of all of them men came to the wedding, that is, believed
in Christ (\work{Quaestiones evangeliorum} I.31). The roads are not the peoples
but what they held, and his answer says nothing of who were first invited or
why they refused, where the other four begin. On the Alleluia's psalm, having
noted that other Latin copies read \latin{Evangelizate}, preach the Gospel, he
asks of its command, ``Unto whom is this addressed, save unto the Evangelists in
prophecy?'' (\work{Enarrationes in Psalmos} 104, 1). Set before the Gospel, the
Alleluia's command to declare God's deeds among the Gentiles and the king's
\latin{Ite ergo ad exitus viarum} are one command in two forms, the psalm's
prophecy and the parable's story. Augustine takes both steps, in two separate
works.

St Gregory the Great gives the highways a different answer, not a different
detail. Preaching to the people of Rome, he names neither Israel nor the
nations in the whole homily. His invited are \latin{amici}, friends, called by
the prophets and then by the apostles, and the ends of the roads are not the
Gentiles: \latin{exitus viarum intelligimus defectus actionum, quia illi
plerumque facile ad Deum veniunt, quos in terrenis actibus prospera nulla
comitantur}, we understand the ends of the ways as the failure of our doings,
because those whom no prosperity attends in their earthly affairs often come
readily to God (\work{Homiliae in Evangelia} 38, 3, 6). Rabanus Maurus's catena
on St Matthew already set Augustine's sentence and Gregory's side by side in
the ninth century, and headed Gregory's as the moral sense. Gregory's answer is
a moral sense of the call, and it stands beside the answer of Irenaeus, Hilary,
Jerome and Chrysostom, for whom the highways are the nations, and beside
Augustine's own, for whom they are the teachings of the Gentiles.

The witnesses divide again over the burned city. Chrysostom reads it of ``the
things that took place under Vespasian and Titus'', which came forty years
after the crucifixion so that God might show his long-suffering. Jerome keeps
the Romans as one of two possibilities beside avenging angels, \latin{de quibus
in Psalmis scribitur: Immissiones per angelos pessimos}, and the psalm he cites
is Ps~77, the Introit's psalm, at its forty-ninth verse. Gregory reads the city
as the flesh of the persecutors, burned with their souls in the eternal fire,
and the armies as the angels (38, 5); Hilary, too, reads the armies as the
armies of heaven and the fire as eternal. The historical and the eschatological
answers stand side by side. St Augustine sets the limit to any wider judgement
inside the Introit's psalm: ``there were even in that people they that
understood, having the faith which was afterwards revealed''
(\work{Enarrationes} 77, 2). And Chrysostom, who blames the first invited for
refusing, requires the garment of the Gentiles who came in as strictly: the
garment ``is life and practice''.

The Collect has its place in this history. The invited went off \latin{alius in
villam suam, alius vero ad negotiationem suam}. For Gregory, \latin{In villam
quippe ire est labori terreno immoderate incumbere, in negotiationem vero ire
est actionum saecularium lucris inhiare}, to go to the farm is to lie
immoderately upon earthly labour, and to go to the merchandise is to gape after
the profits of worldly business (38, 5); for Hilary the farm is the ambition of
the world and the merchandise the love of money. For Chrysostom the excuses
sound reasonable, and that is their lesson: ``though the things which
hinder us be necessary, to set the things spiritual at a higher price than
all''. The Collect asks to be \latin{expediti}, disencumbered in mind and body,
for \latin{quae tua sunt}: free of exactly that weight.

\subsection{The feast that now is, and the feast to come}

The third reading asks where the feast is now. St Gregory opens his homily by
asking whether St Luke's great supper and St Matthew's wedding are one lesson,
and answers that Luke's signifies the eternal and last banquet and Matthew's
the present Church, because at Matthew's wedding one guest who came in is put
out again, and at the eternal banquet no one who has once come in goes out
(38, 1). In his words, \latin{per has regis nuptias praesens Ecclesia designatur,
in qua cum bonis et mali conveniunt}: by this king's wedding the present Church is
signified, in which the bad come together with the good (38, 7). St Augustine
begins his sermon with the same distinction: ``the Holy Scriptures teach us that
there are two feasts of the Lord; one to which the good and evil come, the other
to which the evil come not'' (\work{Sermo} 90, 1). The feast of this Gospel is
the first, its guests are already at the Lord's Table ``which is in this life''
(90, 5), and the whole sermon is gathered into one sentence: ``Be that feast
which now is, received worthily, that we may attain to the other'' (90, 4).

Both preachers turn the parable on the people in front of them. Gregory:
\latin{Sed quia jam, largiente Domino, nuptiarum domum, id est sanctam Ecclesiam
intrastis, solerter, fratres, aspicite, ne aliquid de mentis vestrae habitu rex
ingrediens reprehendat}, since by the Lord's gift you have already entered the
house of the wedding, that is, holy Church, look carefully, brethren, lest the
king coming in find fault with anything in the dress of your mind (38, 9).
Chrysostom, closing his homily on the parable: ``Hear ye, as many as having
partaken of the mysteries, and having been present at the marriage, clothe your
souls with filthy deeds. Hear whence ye were called. From the highway.'' The
three do not name the feast alike. Augustine's is the Lord's Table and ``the
Feast of the Holy Scriptures'' (90, 9); Gregory's is the banquet of the Word,
\latin{qui Scripturae sacrae epulis pascimur} (38, 11), and he does not name the
altar at all; Chrysostom's is the marriage at which his hearers have partaken of
the mysteries. None of them says in so many words that the wedding feast of
this Gospel is the Eucharist. What they share is the time: the feast is now, its
guests are in the Church, and the king has not yet come in. The third reading
hears the chants and prayers of this Mass as the voice of that assembly at that
table, the Church's whole feast of Word and sacrament, of which the three
preachers each name a part.

Two witnesses place the wedding later. St Hilary makes it \latin{vitae coelestis
et in resurrectione suscipiendae aeternae gloriae sacramentum}, the mystery of
the heavenly life and of the eternal glory to be received in the resurrection
(22, 3), and in the second reading the hall filled from the highways is the
Church of the nations awaiting that announcement. Bl.\ Ildefonso Schuster,
commenting on this Mass as a Mass, hears in its Gradual the Christian's evening
oblation and in its Secret and Postcommunion the Eucharist, but he places the
Gospel's banquet in heaven: ``The primary aim of the predestination of souls to
the heavenly banquet'', he writes of the parable, ``is the supreme glorification
of Christ as first-born of the human family and head of the Church'' (\work{The
Sacramentary} III, p.~173). His banquet is the one Augustine sets beside the
present feast as the other, to which the evil do not come. It is Augustine and
Gregory, not Schuster, who find the Church now in this Gospel's wedding. The
witnesses differ over what the wedding signifies, and they do not conflict in
doctrine: Augustine's present feast is received worthily so as to reach the
other, and Hilary places the king's entry in the assembly of the risen, toward
which the present feast looks.

The table is not the garment. Augustine, who calls the present feast the Lord's
Table, denies in the same sermon that the garment is ``the Altar, or That which
is received at the Altar. But no; we see that many eat, and `eat and drink
judgment to themselves'\,'' (90, 5); and in his other sermon on the parable,
``The Sacraments of the Altar the good and bad receive together'' (\work{Sermo}
95, 7). The feast now includes the table, and the garment is what the guest
brings to it. The Secret offers the gifts \latin{oculis tuae maiestatis} and asks that they be \latin{salutaria nobis},
and Schuster reads it of the Eucharist and of those who share in it: since the
sacrament's effect in each person answers to his dispositions, ``we here ask for
such dispositions as may enable the mystical sacramental oblation to exert
without hindrance its full power in our souls'' (p.~173). The Secret's
\latin{salutaria nobis} is Augustine's ``received worthily'' turned into
petition.

\subsection{What the king looks for}

The first reading hears the Gospel from its end, where a man who was called,
came in and sat down is found without the one thing the king looks for. St
Gregory begins by setting aside two answers that lie close to hand:
\latin{Si enim vestem nuptialem baptisma vel fidem dicimus, quis sine baptismate
et fide has nuptias intravit?} If we call the wedding garment baptism or faith,
who has entered this wedding without them? \latin{Quid ergo debemus intelligere
nuptialem vestem, nisi charitatem?} And his reason turns from the guest to the
bridegroom: \latin{Recte enim charitas nuptialis vestis vocatur, quia hanc in se
conditor noster habuit, dum ad sociandae sibi Ecclesiae nuptias venit}, because
our Maker himself had it when he came to the wedding at which he joined the
Church to himself (38, 9). The garment is woven on two beams, the love of God
and of neighbour (38, 10), and tested by the hardest case: \latin{Charitas
autem vera est cum et in Deo diligitur amicus, et propter Deum diligitur
inimicus} (38, 11). The king's greeting is read by it too: \latin{Amice, et non
amice; amice per fidem, sed non amice per operationem}, friend by faith, but no
friend by what he has done (38, 12).

St Augustine, in a sermon against the Donatists, makes the same elimination at
greater length, because the rite that brings good and bad into the hall is
common to both. Not baptism, not the altar, not fasting, not coming to church,
not miracles, since the wicked do all these too; the garment is what the good
have and the bad have not, ``charity out of a pure heart, and of a good
conscience, and of faith unfeigned'' (1 Tim 1:5; \work{Sermo} 90, 5--6). It is
worn ``in the heart, not on the body; for had it been put on externally, it
could not have been concealed even from the servants'' (90, 4), and it grows:
``Let charity be born in thee, if it be yet unborn, and if it be born, be it
nourished, fostered, increased'' (90, 6).

St Jerome names the garment differently, and in doing so joins the parable to
the Epistle's language himself: \latin{Vestis autem nuptialis, praecepta sunt
Domini, et opera quae complentur ex lege et Evangelio, novique hominis efficiunt
vestimentum}, the Lord's precepts and the works fulfilled from the law and the
Gospel, which make the garment of the new man. The guest found without it wears
\latin{veteris hominis exuvias}, the cast skin of the old man (\work{In
Matthaeum} III, PL 26, cols.~160--161), the old man the Epistle's readers were to
put off. Commenting on the Epistle, Jerome says who the new man is, reading the
command beside \latin{Induite vos Christum Jesum} (Rom 13:14): \latin{Iste
quippe est novus homo, quo universi credentes debemus indui atque vestiri}
(\work{In Epistolam ad Ephesios} II, col.~508). St Thomas Aquinas reaches the
same identification in two separate commentaries: on the Epistle,
\latin{principium primum novitatis et renovationis Christus est}, and on the
parable, that the garment is Christ, put on by good work, a holy way of life and
true charity (\work{In Epistolam ad Ephesios} c.~4, lect.~7; \work{Super
Matthaeum} c.~22). Schuster says of this Mass's Epistle, ``This garment is Jesus
Christ himself'' (p.~172). The answers are compatible and not identical:
Gregory and Augustine say what the garment is, Jerome says what it is made of
and whose it is.

The Epistle names what the new man does in four acts, and its commentators read
them as the acts of the new man against the old. Chrysostom takes the first from its reason, ``for we are
members one of another'': ``Let not the eye, saith he, lie to the foot, nor the
foot to the eye.'' On the second: ``let not the sun depart, and leave you both
at enmity.'' And he sums up: ``Seest thou what are the members of the old man?
Falsehood, revenge, theft'' (\work{Homilies on Ephesians} 14). Aquinas allows
anger its first motion and not its delay, \latin{non permittitur mora}
(c.~4, lect.~8--9). Jerome is sharpest on the last: whoever works only so as not to be in
want himself, \latin{et a caeteris contrahit manum}, and draws back his hand
from others, has not kept the Apostle's precept, however he may applaud himself
(on Eph 4:28, col.~512).

The Gospel's last image falls on the same hands. The king orders the man bound
hand and foot, and Gregory reads the binding as already begun: \latin{manus quae
nihil indigentibus tribuunt a bono opere jam ex voluntate ligatae sunt}, hands
that give nothing to the needy are already bound from good work by their own
will (38, 13). Augustine gives the demand its positive form: ``Clothe others,
and ye are clothed yourselves. It is winter, clothe the naked. Christ is naked;
and He will give you that `wedding garment' whosoever have it not''
(\work{Sermo} 95, 7). The thief of the Epistle who works with his hands so as to
have something to give is weaving the wedding garment; the guest whose hands
gave nothing arrives at the feast without it.

\subsection{Gift or work, and the Communion's answer}

If entry is by the king's free call, and the guest is condemned for lacking a
garment, the call seems to demand what it did not give. The witnesses answer in
different ways, and the differences are real. St Hilary makes the garment a
gift: \latin{Vestitus autem nuptialis est gloria Spiritus sancti, et candor
habitus coelestis}, the glory of the Holy Spirit and the brightness of a
heavenly dress, received in the confession of a good answer at baptism and to be
kept \latin{usque in coetum regni coelorum immaculatus et integer} (22, 7). St
Irenaeus says that ``we ought, after our calling, to be also adorned with works
of righteousness, so that the Spirit of God may rest upon us; for this is the
wedding garment'' (IV.36.6). Gregory and Augustine deny that the garment is
baptism, because the unworthy guest is baptized too. Chrysostom holds both sides
in one sentence: ``although to be called and to be cleansed was of grace, yet,
when called and clothed in clean garments, to continue keeping them so, this is
of the diligence of them that are called'' (\work{Homilies on Matthew} 69).

The positions are narrower than they look. Hilary's garment is received in
baptism and must be kept spotless; Gregory's and Augustine's is charity, which a
baptized man may fail ever to put on. They agree that the guest is baptized and
that what he lacks is something a baptized man can lose or never wear, and they
differ over what to call it and where it comes from. Hilary raises a second
difficulty himself: how could a man swept in from the highways be expected to
have a garment? His answer is that the garment is inward, and that God alone
found the unworthy guest; Augustine answers the same way from the servants, who
did not notice.

The Communion gives the formulary's own answer. \latin{Tu mandasti mandata tua
custodiri nimis: utinam dirigantur viae meae}: the Lord's precepts, which Jerome
made the stuff of the garment, are the Communion's \latin{mandata}, and three
expositors of the psalm read its wish as a prayer for grace. St Augustine:
``When thou hearest, `O that,' recognise the words of one wishing; and having
recognised the expression of a wish, lay aside the pride of presumption. \dots\
Therefore if man desireth what God chargeth, God must be prayed to grant Himself
what He enjoineth'' (\work{Enarrationes} 118, Aleph 5). St Hilary, on the same
verse: \latin{nisi a Deo dirigamur, infirmes per naturam nostram erimus.
Adjuvandi igitur per gratiam ejus, dirigendique sumus}, unless we are directed
by God we shall be weak by our nature; we must be helped and directed by his
grace (\work{Tractatus super Psalmos} 118, Aleph 12). And the Greek
\work{Expositiones in Psalmos} printed under the name of St Athanasius say that
the psalmist, knowing that no one can keep the law without help from above,
secures its keeping for himself by prayer. Hilary, who makes the garment a gift
of the Spirit on the Gospel, says on the psalm that the commandments are kept
only as God's grace directs. The garment is God's gift, and it must be put on
and kept; it is kept by praying for what is commanded. The Postcommunion, the
formulary's last prayer, asks for that help, that God's healing work \latin{tuis
semper faciat inhaerere mandatis}. St Robert Bellarmine hears the Communion's
first verse as the lawgiver's authority, a command given ``not by way of advice,
but by strict precept'' (on Ps 118:4), and Schuster hears joy in its
\latin{utinam}.

Each reading gives the Communion its own place. In the first, it is the prayer
by which the garment is put on: what God commands is asked of God. In the third,
sung after the Communion itself, it is the prayer of those who have been at the
table and do not take the direction of their ways for granted; its
\latin{dirigantur} answers the Gradual's \latin{dirigatur}, so that the prayer
that was to rise straight as incense before the sacrifice becomes the prayer
that the ways of those who have eaten be set straight. In the second, it is
where the history ends. The Alleluia's psalm gives as the reason for all it
recites ``That they might observe his justifications'', and St Augustine takes
that verse as the key to the whole psalm: the goods of the land were given so
that God's servants might be free for what is eternal (\work{Enarrationes} 104,
34). The Communion asks to keep those same \latin{iustificationes}. The call
ends not in having been called, but in keeping what the caller commands.

\subsection{The evening sacrifice}

The Gradual sings ``the lifting up of my hands, as evening sacrifice'', and the
witnesses give two answers to what that sacrifice is. St Augustine hears in the
psalm the voice of the whole Christ, which must go on praying ``if tribulation
remain for the Church, for the Body of Christ, even to the end of the world''
(\work{Enarrationes} 140, 2). On the Gradual's verse he turns to the Head: ``For
when the day was now sinking towards evening, the Lord upon the Cross `laid down
His life to take it again'\,''. The members are not left out, ``Still we too are
figured there. For what of Him hung upon the tree, save what He took of us?'',
and he names the sacrifice: ``That then is the `evening sacrifice,' the Passion
of the Lord, the Cross of the Lord, the offering of a salutary Victim, the whole
burnt offering acceptable to God. That `evening sacrifice' produced, in His
Resurrection, a morning offering. Prayer then, purely directed from a faithful
heart, riseth like incense from a hallowed altar'' (140, 3). St Robert
Bellarmine offers the same figure as one possibility among others, ``perhaps,
because the evening sacrifice was of more value as being a figure of the
sacrifice on the cross, that occurred in the evening'', and names the one
through whom such prayer rises: ``Christ is the high priest, for it is through
him, as being our advocate, that we must always pray'' (on Ps 140:2).

St Hilary answers otherwise. To spread out the hands is the posture of prayer;
to lift them shows \latin{excelsi operis officium}, the office of a lofty work,
and when St Paul bids men pray ``lifting up pure hands, without anger and
contention'' (1 Tim 2:8) he \latin{non habitum orandi constituit, sed
praescriptum divinae operationis adjecit}, did not fix a posture of prayer but
added a prescription for divine work (\work{Tractatus super Psalmos} 140, 3).
Asked which evening sacrifice pleases God, Hilary answers from the last
judgement,
``For I was hungry, and you gave me to eat'' (Mt 25:35): \latin{Manuum ergo
elevatio, sacrificium vespertinum est. Non enim sanguine et holocaustis nos, in
quos consummatio saeculorum devenit, sacrificamus Deo} (140, 4). The evening is
the last age, and the lifted hands feed, clothe, welcome and visit Christ in the
least of his brethren. St John Chrysostom, expounding the psalm in Greek, forbids
prayer against enemies and requires that hands raised in prayer be raised
without anger, citing the same verse to Timothy; he adds that the Fathers
appointed the psalm for every evening not for its one phrase about the evening
sacrifice, but as a saving medicine and a cleansing of the sins of the day
(\work{Expositio in Psalmum} 140).

The two answers are not opposed in doctrine: Augustine reads the verse first of
the Head, and Hilary of the hands of the members. But they are two answers to
what the verse's sacrifice is, and the readings divide by them. The first takes
Hilary's. Its lifted hands are the Epistle's working hands, raised in prayer
without the anger the Epistle forbids to outlast the day, and they are the
opposite of the hands Gregory finds already bound, because they have given. The
third takes Augustine's: the evening sacrifice is Christ's, the prayer that
rises with it is the Church's, and the altar from which it rises is the hallowed
altar of a faithful heart. The second hears in the psalm the prayer of a people
still on the road, taking from Augustine the Church that must cry to the end of
the world and from Hilary the people of the last age. Schuster, reading this
Mass's Gradual, calls every Christian offering evening: ``In this present life our oblation to God is
always an evening sacrifice, for it is enveloped in the twilight of faith''
(p.~172).

The Offertory, at which the Gradual's verse returns over the altar in a Mass
with incense, takes up the Introit's word, \latin{Si ambulavero in medio
tribulationis, vivificabis me}, and Augustine's exposition of the verse reads like a gloss on the Introit's
promise: ``whatsoever tribulation thou art in, confess, call on Him; He freeth
thee, He reviveth thee''. He widens the tribulation to the whole of this life,
``whatever prosperity it shine with'', and makes it the condition of being
revived at all: ``otherwise Thou wilt not revive me, unless I walk in the midst
of tribulation'' (\work{Enarrationes} 137, 10). Chrysostom notices that the
psalmist does not say God will drive the tribulation away, but that he will keep
him alive while he remains in the midst of its dangers (\work{Expositio in
Psalmum} 137). Augustine read the verse with the Vulgate's past, and so does
Chrysostom's Greek text, though he reports two other Greek versions with the
future; in the Missal's chant the whole verse looks forward. For the first
reading the walk in tribulation is the road on which the garment is kept; for
the second, the road the servants took and the Church still takes; for the
third, the condition of an assembly that brings its gifts to the altar while it
is still walking in tribulation.

\subsection{Many called, few chosen}

The three readings share their ground. In each, one King calls by grace and
judges those he has called; the Church is a mixed company until he comes in; and
the man without the garment has come in under the call. They differ in the
question they put to the texts, and so the Gospel's last sentence falls on
something different in each. ``For many are called, but few are chosen.''

In the first reading it falls on the guest who came in. Gregory: \latin{Quia
vocati sumus, novimus; si sumus electi, nescimus}, we know that we are called; we
do not know whether we are chosen (38, 14). Jerome reads the king's entry as the
judgement, when the guest found \latin{sub nomine Christiano} without the
garment is rebuked, and says of the vineyard, the building of the house and the
wedding feast alike that what is sought is \latin{non initia, sed finis}, not
the beginning but the end (on Mt 22:14). In the second it falls on the whole
history of the call. Irenaeus: ``For as in the former covenant, `with many of
them was He not well pleased;' so also is it the case here, that `many are
called, but few chosen'\,'' (IV.36.6). Hilary gives the reason:
\latin{Vocatio quidem bonos efficere debuerat, quia sancta est}, the call ought
to have made them good, because it is holy; \latin{sed per vitium inemendatae
voluntatis discrimen est vocatorum}, but by the fault of an unamended will there
is a difference among the called. And he turns the saying: \latin{Non est igitur
paucitas in invitatis, sed raritas in electis}, the few are not few among the
invited but rare among the chosen (22, 6--7). In the third it falls on the feast
that now is, which Augustine tells his hearers to receive worthily so that they
may attain to the other, ``where no bad men shall come''.

Each reading also accounts for something the others leave aside. The first
alone gives full weight to the Epistle, reading its four acts as the garment's
visible side and its new man as Christ put on; with Jerome and Aquinas it hears
Christ named in both lessons. The second alone gives the Introit's psalm and the
Alleluia's command their own voice, and alone takes up the first half of the
parable. The third alone reads the Mass's prayers as the prayers of the
assembly at the altar, from the Gradual and the Offertory to the
Postcommunion's healing work, which Schuster reads of the Eucharist: ``we make our petition that the healing grace
of the Eucharist may mercifully deliver us from our perversity and make us ever
to be devoted to the divine will'' (p.~174).

The anagogical end is the same king's coming in, seen three ways: the judgement
at which the garment is sought (Jerome), a garment that for Hilary is to be kept
spotless until the assembly of the kingdom of heaven; the glory of that kingdom
announced as a wedding when the multitude pleasing to God is gathered (Hilary);
and the other feast, reached from the feast that now is (Augustine), as the
evening sacrifice of the Passion produced in the Resurrection a morning offering
(Augustine on Ps~140). One sentence, one judgement, three times: the guest's,
the history's, and the table's.
